% Equivalencias de las órdenes de tcee.sty para pandoc (HTML y Word).
% construir-tema.py lo antepone al LaTeX exportado de LyX. Si se añade una
% orden nueva a tcee.sty, hay que añadir aquí su equivalente.
\newcommand{\autor}[1]{\textsc{#1}}
\newcommand{\concepto}[1]{\textbf{#1}}
\newcommand{\vertema}[1]{\textcolor{tceegris}{[ver tema #1]}}
\newcommand{\vertemacorto}[1]{\textcolor{tceegris}{#1}}
\newcommand{\fuente}[1]{\textit{Fuente: #1}}
\newcommand{\grafico}[2][0.75\textwidth]{\includegraphics[width=#1]{#2}}
\newcommand{\tceetitulo}{}
\newcommand{\bibliografia}{}
\newcommand{\tcetema}[3][]{}
\newcommand{\tcefecha}[1]{}
\newcommand{\minutos}[1]{\textcolor{tceeminutos}{#1}}
\newcommand{\vect}[1]{\vec{#1}}
% \tiempo{2,5} → 2′30″ (lo calcula tcee.lua, como \minutos)
\newcommand{\tiempo}[1]{\textcolor{tceetiempo}{#1}}
% Gráfico pequeño en una nota al pie (tcee.lua y tema.css lo centran)
\newcommand{\graficonota}[2][0.45]{\begin{notagrafico}\includegraphics[width=#1\linewidth]{#2}\end{notagrafico}}
% Fin de demostración (tcee.lua lo pone a la derecha en la web; fuera de las fórmulas)
\newcommand{\findemo}{\textcolor{tceefindemo}{$\blacksquare$}}
% \funcionclave y \optim los expande construir-tema.py antes de pandoc
% figuracaja y tablacaja (gráfico o tabla dentro de una caja) los convierte
% construir-tema.py en figure y table antes de pandoc
